\documentclass[12pt]{article}
\usepackage[utf8]{inputenc}
\usepackage{amsmath, amssymb, amsfonts}
\usepackage{geometry}
\usepackage{booktabs}
\usepackage{enumitem}
\geometry{margin=1in}
\setlist[enumerate]{leftmargin=*}

\title{Practice Problem Set 1\\ \vspace{15pt} \large ECON 101 -- Macroeconomics}
\author{}
\date{Fall 2025}

\begin{document}
\maketitle



\begin{enumerate}
  \item 
  Explain why scarcity is considered the fundamental problem of economics. Provide two real-world examples of trade-offs individuals face, and discuss how opportunity cost is involved.

  \item 
  Why do economists use models and assumptions when studying real-world issues? Briefly discuss the strengths and limitations of using models in macroeconomics.

  \item
  Give one example each of a positive statement and a normative statement related to government spending. Explain why the distinction matters in economics.

  \item
  Using the concept of comparative advantage, explain how trade can make both trading partners better off, even if one is more productive in producing all goods.

  \item
  Explain what a point inside the PPF, on the PPF, and outside the PPF represent. Give one example for each.
\end{enumerate}



\begin{enumerate}[resume]
  \item 
  A student can either write an essay (which takes 5 hours and earns her a grade improvement worth the equivalent of \$80) or work part-time at a café (which pays \$15/hour).  
  \begin{enumerate}
      \item What is the opportunity cost of writing the essay?  
      \item Should the student write the essay or work? Explain using economic reasoning.
  \end{enumerate}

  \item 
  Consider an economy that produces only two goods: wheat and cars. The following table shows maximum output combinations:

  \begin{center}
  \begin{tabular}{c c c}
    \toprule
    Combination & Wheat (tons) & Cars (units) \\
    \midrule
    A & 0   & 100 \\
    B & 200 & 90  \\
    C & 400 & 70  \\
    D & 600 & 40  \\
    E & 800 & 0   \\
    \bottomrule
  \end{tabular}
  \end{center}
\newpage
  \begin{enumerate}
      \item Draw the PPF.  
      \item Calculate the opportunity cost of moving from combination B to C.  
      \item Is the opportunity cost constant, increasing, or decreasing? Explain.  
  \end{enumerate}

  \item
  Two countries, Alpha and Beta, can each produce cloth and wine. Their daily production possibilities are:

  \begin{center}
  \begin{tabular}{c c c}
    \toprule
    Country & Cloth (units) & Wine (units) \\
    \midrule
    Alpha & 60 & 30 \\
    Beta  & 40 & 20 \\
    \bottomrule
  \end{tabular}
  \end{center}

  \begin{enumerate}
      \item Calculate the opportunity cost of cloth and wine in each country.  
      \item Identify which country has comparative advantage in each good.  
      \item Explain how both countries can gain from trade.  
  \end{enumerate}



\end{enumerate}

\end{document}
